% 07-cahier-des-charges.tex — Réalisation 1/4 : exigences pédagogiques et
% techniques, puis méthodologie de développement assistée par IA (≈ 8 pages,
% dont la méthodologie déjà rédigée dans sections/methodologie.tex).

\chapter{Cahier des charges et méthode de travail}
\label{chap:cahier}

\section{Exigences pédagogiques}

\begin{arediger}{1 page}
Traduire l'état de l'art en exigences vérifiables : chaque page montre le
phénomène, \emph{toutes} les forces avec un code couleur constant (poids
rouge, réaction verte, tension violette, vitesse en tirets, accélération en
pointillés, résultante épaisse), un diagramme des forces toujours visible,
les formules sous trois formes (symbolique, valeurs substituées, résultat
avec unité), les hypothèses écrites, des graphes temporels synchronisés et
un contrôle complet du temps. Notations du formulaire du cours.
\end{arediger}

\section{Exigences techniques}

\begin{arediger}{1 page}
Site statique (HTML, CSS, JavaScript sans bibliothèque) qui fonctionne par
double-clic (\fichier{file://}), sur PC, tablette et téléphone ; physique
isolée dans un fichier par exercice, facile à auditer ; mode jeu supprimable
en quatre étapes. Expliquer pourquoi chaque contrainte a été choisie.
\end{arediger}

\input{sections/methodologie}
